\documentclass[11pt]{article}

% ============================================================================
% Grant application (draft). Bracketed [ ... ] items are placeholders to fill
% in or confirm before submission: co-PI name/affiliation, sponsoring ministry,
% figures, dates, and the exact Masclans et al. citation.
% Compiles with a standard TeX distribution: pdflatex kazakhstan-innovation-os.tex
% ============================================================================

\usepackage[margin=1in]{geometry}
\usepackage{booktabs}
\usepackage{enumitem}
\usepackage{titlesec}
\usepackage[hidelinks]{hyperref}
\usepackage{microtype}

\titleformat{\section}{\large\bfseries}{\thesection}{0.6em}{}
\titleformat{\subsection}{\normalsize\bfseries}{\thesubsection}{0.6em}{}
\setlist{itemsep=2pt,topsep=3pt}
\setlength{\parskip}{4pt}
\setlength{\parindent}{0pt}

\title{\bfseries A National Innovation Operating System:\\
Converting Scientific and Human Capital into Economic Value in the Republic of Kazakhstan}

\author{
  Sharique Hasan\\
  \small The Fuqua School of Business, Duke University\\[4pt]
  [Co-Principal Investigator], [Ministry / Institution], Republic of Kazakhstan
}

\date{[Month Year]}

\begin{document}
\maketitle

\begin{abstract}
\noindent
Kazakhstan has set the strategic objective of diversifying its economy beyond
extractive industries toward a knowledge-based economy driven by science,
technology, and a modernized workforce. The binding constraint on that
transition is not the absence of scientific or human capital, but the absence of
an \emph{instrument} that can measure it, identify where it is most valuable, and
convert it into economic activity at national scale. This proposal delivers that
instrument. Building on a validated machine-learning approach to identifying
commercially valuable science (Masclans et al.), a deployed research-scoring
platform (\texttt{scientifiq.ai}), a workforce-redesign engine developed under an
OpenAI research grant, and the \texttt{superadditive.app} authoring and delivery
infrastructure, we propose a \emph{National Innovation Operating System} (NIOS):
an integrated set of AI instruments that (i) score the commercial, scientific,
and social potential of the nation's entire research base; (ii) assemble a
fundable, diversified deep-technology portfolio; (iii) map and reconnect the
scientific diaspora; (iv) measure and redesign the AI exposure of the national
workforce; (v) profile the small- and medium-enterprise sector's technology
readiness; and (vi) simulate the marginal returns to public research investment.
Crucially, each component already exists in working form; the grant validates the
system on Kazakhstani data, deploys it on sovereign infrastructure, and stands up
the closed loop from measurement to policy. We therefore propose not a program of
promised research but the national deployment of a functioning system.
\end{abstract}

\section{Motivation and Vision}

Resource-rich economies face a well-documented development trap: the very
abundance that finances the state can crowd out the institutions and incentives
required to build a diversified, innovation-driven economy. Kazakhstan's national
strategy explicitly seeks to escape this trap by cultivating high-value
knowledge industries, commercializing university research, digitizing government
and industry, and equipping the workforce for an economy in which artificial
intelligence is a general-purpose technology. These are the right objectives.
What has been missing is a way to act on them with precision.

A government seeking to diversify confronts a series of concrete questions.
Which fields of its national research base are genuinely commercially valuable,
and which are strong only academically? Which specific inventions, sitting today
in university laboratories, should a sovereign fund back? Where are the country's
scientists---at home and in the diaspora---and which collaborations would create
the most value? Which occupations are most exposed to AI, and how should those
jobs be redesigned rather than merely eliminated? Which firms are ready to adopt
new technology, and which require support? And, given a finite research budget,
where does the next unit of public investment produce the greatest return?

Historically these questions have been answered, if at all, by committee,
anecdote, or imported consultancy. We propose instead to answer them
\emph{computationally and continuously}, using instruments that read the
nation's scientific and economic text at scale and translate it into decisions.
The scientific premise is that the commercial and social potential of research,
and the AI exposure of work, are \emph{predictable from text} with useful
accuracy; the engineering premise is that the instruments to do so already
exist and can be pointed at Kazakhstan. This proposal joins the two.

\section{Intellectual and Technical Foundations}

The proposed system rests on four foundations that the principal investigators
have already built and validated.

\subsection{A validated science of commercial potential (Masclans et al.)}
The intellectual core is a machine-learning approach that predicts the
commercial, scientific, and social potential of scientific work directly from
its text, learned from historical relationships between published research and
downstream outcomes such as patenting, licensing, and firm formation. This body
of work establishes that ``valuable science'' is identifiable at scale and in
advance, rather than only in retrospect---the necessary precondition for any
national commercialization strategy. It supplies both the measures the system
uses and the scholarly validation that those measures are meaningful.

\subsection{A deployed research-scoring platform (\texttt{scientifiq.ai})}
These measures are operationalized in \texttt{scientifiq.ai}, a production
application-programming interface that scores research abstracts along commercial,
scientific, and social dimensions and supports a suite of higher-order
instruments: institutional benchmarking, invention scoring and ranking,
licensing-brief generation, technology-domain scans, collaborator and
co-founder identification, and portfolio optimization. The platform is backed by
frozen scientific-language-model embeddings with calibrated predictive heads and
is served in production; it is, in effect, the measurement layer of the proposed
system, already engineered and running.

\subsection{A workforce-redesign engine (OpenAI research grant)}
Under an OpenAI research grant, the investigators developed instruments that
measure the AI exposure of occupations and guide the \emph{redesign} of work---
distinguishing what only humans can lead, own, and judge from what AI can search,
structure, and translate---and deliver this as a structured, hands-on exercise
rather than a static report. This engine converts the abstract anxiety of
``automation'' into concrete, occupation-level reskilling and job-redesign plans,
and is the human-capital counterpart to the science-commercialization layer.

\subsection{An authoring and delivery infrastructure (\texttt{superadditive.app})}
Finally, \texttt{superadditive.app} provides the delivery and authoring layer: a
platform through which these instruments are configured, run with real
populations (cohorts of professionals, civil servants, or students), and
observed. It already includes multimodal enterprise profiling, live facilitated
activities, and---critical for a government partner---the ability to route all
computation through an organization's \emph{own} private models and
infrastructure (Section~\ref{sec:sovereignty}).

Individually, these are research artifacts and products. Integrated and pointed
at a single nation, they become an operating system for the innovation economy.

\section{The National Innovation Operating System}

We organize the work into six instruments (work packages, WP1--WP6). Each is
built on the foundations above and already exists in working form; the grant's
contribution is validation on Kazakhstani data, sovereign deployment, and
integration into a closed decision loop. Together they trace the value chain from
scientific knowledge to economic activity: \emph{measure} the science,
\emph{select} the bets, \emph{connect} the talent, \emph{redesign} the work,
\emph{map} the firms, and \emph{allocate} the funds.

\subsection*{WP1 --- The Kazakhstan Science-to-Value Atlas}
We score the nation's entire research output---every paper, researcher,
institution, and field for which text is available through open bibliographic
sources---along the commercial, scientific, and social dimensions. The result is
a living atlas of where Kazakhstan's commercializable science actually resides:
its genuine strengths, its overlooked commercial assets, and its gaps relative to
comparator states. This transforms research-policy debate from opinion to
evidence and gives ministries a defensible basis for concentrating support.

\subsection*{WP2 --- The Sovereign Deep-Technology Portfolio}
Applying invention scoring, ranking, and portfolio optimization, we assemble
from the national research base a \emph{diversified, fundable} portfolio of
high-potential inventions, each accompanied by an automatically generated
licensing brief situating it against comparable science and the surrounding
patent landscape. The deliverable is a decision-ready shortlist for a sovereign
technology fund or a national tech-transfer office---exactly the artifact such
bodies currently lack.

\subsection*{WP3 --- The Diaspora and Collaboration Engine}
We map Kazakhstani-origin scientists worldwide and identify the collaborations
and co-founder pairings that would create the most value, given the national
research profile. This directly addresses brain drain: rather than lamenting the
diaspora, the system provides a concrete plan to reconnect it and to seed
science-based ventures.

\subsection*{WP4 --- National AI Exposure and Job-Redesign Program}
Beginning with the civil service and two priority sectors, we measure the AI
exposure of occupations and deliver, at scale, structured job-redesign and
reskilling plans. The output is a workforce-modernization blueprint grounded in
the specific tasks Kazakhstani workers actually perform, positioning AI as a tool
for augmenting the labor force rather than displacing it.

\subsection*{WP5 --- The Enterprise Innovation Census}
Using multimodal profiling, we build a living panel of the small- and
medium-enterprise sector's technology readiness and needs. The panel supports
targeting of support programs and, because it is longitudinal, the measurement of
technology diffusion over time---an evaluation capability most innovation
programs never obtain.

\subsection*{WP6 --- The Research-Allocation Simulator}
Finally, using value-to-go estimation, we simulate the marginal commercial and
social return of public research investment across fields, giving policymakers a
principled answer to the question of where the next unit of funding should go,
and what reachable ceiling it can be expected to approach. This work package is
the system's research frontier and its most novel scholarly contribution.

\section{Technical Approach}

Methodologically, the system is built from text. Scientific and occupational
corpora are embedded using scientific-language models; calibrated predictive
heads map these embeddings to the commercial, scientific, and social potential
measures validated in the foundational work. Higher-order instruments compose
these scores: invention ranking orders disclosures by predicted potential;
portfolio optimization selects a set that maximizes aggregate potential subject
to a diversity constraint, so that the recommended portfolio is not merely a list
of top-scoring items but a genuinely spread set of bets; and value-to-go
estimation, trained through self-play over sequences of research and
commercialization decisions, estimates not just an item's current value but the
value reachable from it---distinguishing already-mature work from work with
headroom. The workforce instruments apply the same text-first logic to
occupational descriptions and to the artifacts workers produce.

Two properties make the approach suitable for national deployment. First, it is
\emph{automatic and continuous}: once pointed at a corpus, it scores it without
manual review, and re-scores as the corpus grows. Second, it is
\emph{interpretable at the decision level}: every recommendation is accompanied
by the evidence---the comparable science, the patent landscape, the specific
tasks---that a human decision-maker needs in order to act and to be accountable.

\section{Data Sovereignty and Governance}
\label{sec:sovereignty}

For a government partner, data sovereignty is not a feature but a precondition.
The delivery infrastructure supports per-organization configuration of the
underlying AI provider: Kazakhstan can route all computation through its own
private or self-hosted models, running on national infrastructure, so that
sensitive research, enterprise, and workforce data never leave the country or
reach a third-party model. The system defaults to this sovereign configuration
where the partner requires it and falls back to shared infrastructure only for
non-sensitive tasks and only by explicit choice. All personal and
enterprise-level data are governed under an agreed data-protection framework,
processed for the stated public-policy purposes alone, and retained under the
partner's control. This architecture allows the system to deliver national-scale
intelligence while keeping the nation's data within the nation's authority.

\section{Work Plan and Timeline}

The work proceeds in three phases over [24] months. Because the instruments
already exist, the early phase emphasizes data integration and validation rather
than construction, allowing decision-ready deliverables within the first two
quarters.

\begin{center}
\begin{tabular}{@{}p{2.4cm}p{5.2cm}p{6.2cm}@{}}
\toprule
\textbf{Phase} & \textbf{Period} & \textbf{Milestones} \\
\midrule
Phase I: Foundation & Months 1--6 &
Data-sharing and sovereignty agreements; ingestion of the national research
corpus and institution list; sovereign-model deployment; first Science-to-Value
Atlas (WP1) and initial deep-technology portfolio (WP2). \\[4pt]
Phase II: Deployment & Months 7--16 &
Diaspora and collaboration engine (WP3); AI-exposure measurement and first
job-redesign cohorts in the civil service and two priority sectors (WP4);
enterprise census pilot (WP5). \\[4pt]
Phase III: Integration & Months 17--[24] &
Research-allocation simulator (WP6); integration of the six instruments into a
single decision loop; validation study, policy handbook, and capacity transfer to
the partner institution. \\
\bottomrule
\end{tabular}
\end{center}

\section{Expected Outcomes and Policy Impact}

The tangible outputs are decision instruments in daily use by the partner
ministries: an atlas the research-policy body consults, a portfolio the sovereign
fund draws from, a diaspora plan the talent agency executes, a reskilling
blueprint the civil-service and labor bodies implement, an enterprise panel the
SME-support programs target, and an allocation simulator the science-funding body
consults. The intellectual outputs include peer-reviewed validation of the
potential and value-to-go measures on a national corpus, and---through WP6---a
novel contribution to the economics of public research allocation. The strategic
outcome is a repeatable, sovereign capability: an operating system Kazakhstan
owns and can run indefinitely, converting its scientific and human capital into
economic value long after the grant concludes.

\section{Team and Feasibility}

The principal investigator has built the scientific measures, the scoring
platform, the workforce-redesign engine, and the delivery infrastructure on which
this proposal depends; the co-principal investigator provides the governmental
mandate, data access, and institutional pathway to adoption within the Republic of
Kazakhstan. Feasibility is unusually high for a proposal of this ambition
precisely because its central risk---whether the instruments can be built---has
already been retired. What remains is validation on national data and deployment,
which the team is positioned to execute immediately.

\section{Budget Summary}

We request [amount] over [24] months. Principal categories are: personnel for
data engineering, validation, and deployment; sovereign compute and model-serving
infrastructure; data acquisition and integration; fieldwork for the enterprise
census and workforce cohorts; and dissemination and capacity transfer to the
partner institution. A detailed budget and justification are provided in the
supplementary materials.

\begin{thebibliography}{9}
\bibitem{masclans} Masclans, [Initials], Hasan, S., et al.
  \emph{[Title: machine-learning identification of commercially valuable
  science].} Working paper, [Year]. % TODO: confirm exact citation.
\bibitem{scientifiq} Hasan, S. \emph{scientifiq.ai: a platform for scoring the
  commercial, scientific, and social potential of research.} Technical report,
  [Year].
\bibitem{superadditive} Hasan, S. \emph{superadditive.app: infrastructure for
  authoring and running AI-based learning and assessment.} Technical report,
  [Year].
\bibitem{openai} Hasan, S. et al. \emph{Redesigning work in the age of AI:
  measuring occupational exposure and guiding job redesign.} OpenAI research
  grant, [Year].
\end{thebibliography}

\end{document}
